% LREC 2026 Example; 
% LREC Is now using templates similar to the ACL ones. 
\documentclass[10pt, a4paper]{article}

\usepackage[review]{lrec2026} % this is the new style
% the 'review' option anonymizes the paper following submission guideline
% the 'final' option produces the camera ready version (non anonymized)
% default version is 'final', so use review option for submission

\title{Title of the LREC 2026 Paper (Title in 14-point Bold, Capitalized in Titlecase)}

\name{Author1, Author2, Author3} 

\address{Affiliation1, Affiliation2, Affiliation3 \\
         Address1, Address2, Address3 \\
         author1@xxx.yy, author2@zzz.edu, author3@hhh.com\\
         \{author1, author5, author9\}@abc.org\\}

\abstract{
Each paper must include an abstract of 150 to 200 words in 9 pt with interlinear spacing of 10 pt. The heading Abstract should be centered, 10 pt bold. This short abstract will also be used for producing the Booklet of Abstracts (PDF) containing the abstracts of all papers presented at the Conference.
 \\ \newline \Keywords{keyword1, keyword2, keyword3} }

\begin{document}

\maketitleabstract

\section{Introduction}

Text to speech (TTS) systems are typically built to use phones rather than raw text. Even for languages with shallow orthographies, the relationship between the written and spoken form is often indirect, so phones act as an intermediary that are intended to ensure the correct sounds are produced.

Pronunciation dictionaries are often hard to find, and are never complete in a living language. For this reason, the task of grapheme to phoneme (G2P) conversion exists: to automate the production of a phonetic representation of words not in the dictionary.

Grapheme-to-phoneme (G2P) conversion is typically considered in terms of a single language, but languages do not exist in isolation: words, names in particular, from other languages make frequent appearances.

While monolingual G2P systems perform well within their target language, they struggle to adapt to words borrowed from other languages. Unassimilated loanwords follow the rules of their source language, and can even have a negative impact on the source language if too many are introduced into the training data of a G2P system, by introducing confusion. For instance, even state-of-the-art end-to-end TTS systems such as DiTTo-TTS \cite{lee2025dittotts} occasionally mispronounce words from well-resourced languages like Polish, despite large multilingual datasets.

A further complication arises from how text is marked up for synthesis. Standards such as the Speech Synthesis Markup Language (SSML) \cite{w3c:ssml11} permit the explicit specification of language on multiple levels—paragraph, sentence, or word--but these annotations are only useful if the G2P component itself can interpret them correctly.

% TODO: Revise this paragraph so it states the intended direction from junk.tex, not only the Phonetisaurus baseline.
In this work, we explore the extent to which multilingual or mixed-language G2P training improves—or degrades—performance across languages. Using the Braxen lexicon \cite{tannander-edlund-2025-braxen}, a Swedish TTS resource containing extensive foreign-language entries, we compare monolingual, multilingual, and untagged models to quantify how cross-language exposure affects phoneme error rate (PER). Our goal is to better understand how multilingual data and language tagging influence pronunciation accuracy in practical TTS systems.

% In this work, we attempt to use a combination of data selection and task prompting to guide a (by)T5 model, which has already been shown to perform well at the G2P task, towards being able to produce phonetic transcriptions of words from one language using the phoneset of another with minimal intervention.
% "minimal intervention", yuck. It's G2P, not a UN peacekeeping mission
% TODO: This seems like the intended contribution statement, but it needs revision before becoming running text.

\section{Background \& related work}

Recent work in text to speech (TTS), such as \citet{gao2023e3tts}, attempts to move towards ``end to end'' synthesis, without the use of a phonetic transcription. As impressive as they are, they do not always pronounce everything correctly: on the samples page\footnote{\url{https://ditto-tts.github.io/}} for DiTTo-TTS \cite{lee2025dittotts}, for example, in the second Polish example the word ``niektórzy'' is mispronounced: Polish has quite a shallow orthography and a relatively large amount of publicly available data.

% SSML etc.?
Speech Synthesis Markup Language (SSML) \cite{w3c:ssml11}, the most widely used standard for representing to a TTS engine how a piece of text should be read, allows language to be specified on multiple levels, including paragraph, sentence, token, and word. This degree of control over generated speech is desirable in many contexts, particularly in contexts where the end-users rely on a spoken form, such as those with print difficulties.

One typical option for handling words of foreign origin is to perform a phoneset mapping, where the output of a G2P system for the other language is mapped to that of the current language. This requires a separate G2P component for each language individually, in addition to the mapping, which increases the complexity of the text processing pipeline.

% In a rule-based paradigm, there are two common options for adapting a system designed for one language to the phoneset of another: either by editing the G2P rules themselves so that they output phones appropriate for the second language, or by explicitly mapping the phones it outputs to those of the other language.

% Pivots (or bridge languages) have long been used in machine translation to augment lexicons in rule-based machine translation \cite{forcada2011apertium}, or phrase tables in statistical machine translation \cite{kumar-etal-2007-improving,CHEN08.733}.

% Google's Neural Machine Translation \cite{johnson-etal-2017-googles} learned an implicit pivot through the use of a single set of shared modules with a prepended token to signify target language, the alignments tended to converge in a space akin to an interlingua. In addition to improving performance among less resourced language pairs, they found that ``zero-shot'' translation became possible for language pairs that weren't seen in training, but for which a pivot was available.

% Need to say something about T5, because it's a key step
% That is: the artificial token is a precursor to the kinds of prompts used in T5, but it broadly equivalent to using the same kind of token when concatenating FSTs.

% A non-goal for this work is the kind of ``zero shot'' G2P claimed in \cite{zhu22_interspeech}. Although it does seem, as they claim, to be a useful basis for future G2P models for under-resourced languages, we consider the concrete case of G2P as a component of a text-to-speech system which must, as a result of human limitations, be bounded to the phoneme inventory of the speaker. Our intention, rather, is to prime the G2P system through data- and task-selection to produce something that is within the speaker's range, and hopefully broadly representative of how a speaker of one language pronounces loanwords from another.
% TODO: This non-goal is central to the intended framing, but should be cleaned up and moved into running text.
% TODO: Make the distance from broad ``zero-shot G2P'' explicit: the paper is not trying to produce arbitrary IPA for an unseen language, but to solve a TTS-bounded pronunciation problem where a foreign word is rendered inside the target speaker's usable phoneset.

\section{Method}

% OK cool. Now add placeholders in the method. Minimally two sentences per process step you need to go through to get this done. Language/writing quality does not matter at all though.

\subsection{Data}

\subsubsection{Braxen}

Braxen is an off-shoot of the internal TTS lexicon maintained by the Swedish Agency for Accessible Media (MTM), whose mission includes the production of material for people with print difficulties. This involves creating audiobooks and TTS services for news, so the lexicon is kept up to date with new vocabulary on an on-going basis.

As such, some of the vocabulary in Braxen reflects the realities of text ``in the wild'': it includes a large selection of words from other languages. From our perspective, the detail that is most interesting is that it includes language tags for these foreign words.

%As the primary mission of MTM is not 
As MTM is primarily concerned with the practical matter of producing spoken forms of written material, rather than lexicography for its own sake, a number of errors and inconsistencies appear in the data in relation to the tagging of languages: the languages are tagged using a three-letter language code: in the case of Latin (`lat'), this is a single character away from the character codes of Latvian (`lav') and Lithuanian (`Lit'): consequently, we find a small number of words from those languages tagged as Latin.

By far the most represented foreign language in Braxen is English.
English is taught as a mandatory subject in schools across Sweden \cite{norrby2014english}: according to the 2024 Eurobarometer report, 90\% of Swedes report feeling comfortably conversational in English \cite{eurobarometer2024languages}, and according to the EF English Proficiency Index 2024, Swedes rank 4th in the world in terms of their English skills \cite{efepi2024}.

In addition to being the primary languages of neighbouring countries, Danish, Swedish, and Norwegian form a dialect continuum, and there is some degree of mutual intelligibility between their speakers. As such, Danish and Norwegian words are among the best represented in Braxen. However, as reflects the nature of text in the wild, many of these words are spelled using Swedish letters: the Norwegian and Danish alphabets have the letters `æ' and `ø' where Swedish has `ä' and `ö', respectively, with `æ' additionally rewritten as `ae'.

% In addition to the part of speech information also contained in Braxen, the NST lexicons include information about named entities, which is quite detailed in terms of category and subcategory. It also contains compound information, which is useful to the grapheme-to-phoneme process: the letter sequence `sch' is a trigraph representing a single sound in `dusch' (\textit{shower}), but represents `s' and `ch' in `potatischips' (\textit{potato chips}), for example.

% Compounds in NST are split into their composite parts, including the linker: while it might be more correct in terms of compounds to think of the linking `s' as a separate entity, rather than treating it as part of the noun (usually, making it identical to the genitive form), for the purposes of syllabification, it is preferable to keep this `s' attached to the correct noun: e.g., ``Östermalmstorg'' ought to be syllabified in terms of ``Östermalms'' $+$ ''torg'', not ``Östermalm'' $+$ ``*storg''.

% This particular example demonstrates a limitation with our method: ``Östermalm'' itself is a compound: ``öster'' + ``malm'' (\textit{east suburb}), but this cannot be recovered from the Braxen data, as the transcription lacks a compound marker. This is, however, likely to be preferable for other purposes: none of the compounds that contain it refer to anything other than the place.

% Although the named entity information is less information for our purposes, we choose to maintain it; our conversion of the NST data is complete, though modernised. First, the data has been converted to Unicode; transcriptions, which occupied a fixed number of fields in the original data, have been converted to a list structure; finally, in addition to maintaining the original transcription, we also add an IPA conversion: although the transcriptions are claimed to use SAMPA, in reality the original transcriptions are rather loosely based on SAMPA.
% TODO: Decide whether NST belongs in the same paper, an appendix, or a later extension.

% TODO: Baseline method from lrec2026-example.tex starts here.
As a first pass, we attempted to find misclassified words by checking for characters that are typical of the language each word was tagged as belonging to. For instance, Swedish words rarely contain the letters `ç' or `ñ', whereas these are common in French and Spanish, respectively. We compiled a list of such characteristic letters for each language represented in Braxen and used it to flag cases where the spelling and language tag appeared inconsistent. While this simple heuristic highlighted a number of misclassified words, it generated too much noise: very few of the characters in the languages represented in Braxen are truly unique to that language alone, and even when that is the case, there are too many instances of words with incorrect accents, such as `\`a' instead of `\'a' in Spanish.

We concentrated our efforts by selecting only languages with 1,000 or more entries; as a replacement form of filtering, we passed each candidate word to a spelling checker, selecting only those that were marked as correct. One clear limitation of this approach is for languages whose script is not Latin: the majority of words classified as Russian in Braxen are transliterated, which reduced the initial number of entries from 1,320 to just 7.

Another, somewhat surprising limitation of this approach was the limits of the spell checking dictionaries for certain languages. For Spanish, for example, a large number of proper names we know to be correct were simply absent from the default dictionary.

Of the initial set of languages, only English, Danish, French, Spanish, Italian, and Latin were retained.

For our comparisons, we trained G2P models using Phonetisaurus \cite{NOVAK_MINEMATSU_HIROSE_2016} in three configurations: individual models, one for each source language; a combined ``raw'' model, which mixes all vocabulary, regardless of origin; and a ``merged'' model, where non-Swedish words were given a prefix unique to that language.

In addition to these three configurations, we also tested the results of the merged model without providing the prefix. For each language and configuration, we calculated the phoneme error rate (PER). The results are in table \ref{tab:per_results}: \textbf{WL} refers to the individual models, \textbf{MWL} refers to the merged model tested with a language-specific prefix; \textbf{WLM} refers to the merged model tested without providing the prefix; and \textbf{RAW} refers to the model that combines all entries without any prefix.

% TODO: Prompted G2P method from junk.tex belongs here once it can be written as running text.
%% charsiu g2p just uses a language prefix
%% doing something similar... `g2p|lang=sv` extended to `g2p|lang=sv|from=en` etc.

% \item Start from a pretrained multilingual T5 model
% byT5; a lot of the language around this has been kind of vague. The Charsiu paper, which is the main one to refer to right now, muddies the waters by having some that they pre-trained themselves (I guess as a crappy, easily beatable baseline), and being vague about which ones are fully trained and which are merely fine-tuned.

%\item Introduce a light adaptation step: teach the model to map words to phonemes using simple prompts (g2p|lang=sv: word).

% \begin{itemize}
% \item Use small language-specific G2P datasets for that alignment.
% \end{itemize}

% Basic monolingual G2P is a baseline task, and it's foundational, but I'm not interested in it for its own sake.

% The next step is this:
% \item Introduce specific vocabulary that is additionally tagged with the donor language: \texttt{g2p|lang=sv|from=en}. The equivalent SSML would be this:

% \begin{verbatim}
% <speak xml:lang="sv-SE">
%   Jag ser fram emot <lang xml:lang="en-US">weekend</lang>!
% </speak>
% \end{verbatim}

% Probably for Swedish it would also be beneficial to introduce a compound splitting task, because the 'sch' in dusch is not the same as in potatischips. 
% JE Honestly it makes little (as in "not none at all") sense to do that within the G2P task. It really is better dealt with as preprocessing. For a whole slew of reasons. If usefulness, rather than the ability to use buzzwords (i.e. an empty "end-to-end") is the goal, then baking segmentation into the G2P task is an at least relatively bad idea. Similar to baking language detection into it, mind you...

% TODO: Replace these commented bullets with an implemented experimental description once the notebooks/scripts exist.

% TODO: Pivot task method from junk.tex belongs here once it can be written as running text.
% Hmm... I know you said to skip the research questions, but to put it in text form, the things are:
% Is the prompt going to be enough? In a full text system, there's enough context-based information to learn degrees of synonymity, the broad equivalents of parts of speech, etc. The prompt alone is enough to introduce sufficient context to know that a different set of outputs are expected, but it's not necessarily the case that it will be sufficient to align the 'from=X' to monolingual 'lang=X'; having an explicit second task, 'ipa2ipa|from=X|to=Y' is how I intended to overcome that, an explicit pivot. Testing this is what I'm stuck on: are both needed?

% JOR: the base case I'm looking at, if everything else fails, is of a mechanism that can be plugged into the text processing layer of a TTS system that can be **specifically** told `X is an English word, pronounce it as a Swede would'; what I'm betting on is that giving a model with T5-level expressivity an explicit pivot can cut down on the amount of explicit data needed to achieve _that_ task in an otherwise multilingually trained G2P system.

% \item Add a pivot task (ipa2ipa|from=en|to=sv) to test if explicit bridging improves cross-language consistency.
% JE Ok so "add" is not really right, it's more been the primary goal all along?
% JOR Right. Whether explicitly or via a pivot, a G2P that can directly handle foreign words according to a single language's phoneset is a relatively rare thing, so even if the pivot fails the fallback is still a win.

%FIXME: placement

% TODO: This is the intended direction of the paper. It needs a concrete method description, data construction recipe, and ablation plan.

% TODO: Evaluation plan from junk.tex belongs here once the actual experiment matrix is fixed.
% \item Evaluate:

% \begin{itemize}
% \item Zero-/few-shot G2P accuracy

% \item Cross-language consistency

% \item Evidence of shared phonological representations (embedding similarity)
% \end{itemize}

% \end{itemize}

% TODO: Define the actual evaluation set and metrics for the prompted/pivot experiments.

\section{Results}

% TODO: Baseline results from lrec2026-example.tex start here.
\begin{table*}[ht]
\centering
\caption{Phoneme Error Rates (PER) for individual and merged G2P models. Lower is better.}
\label{tab:per_results}
\begin{tabular}{lcccccccc}
\toprule
\textbf{Language} & \textbf{WL} & \textbf{MWL} & \textbf{WLM} & \textbf{RAW} &
$\Delta_{\text{MWL-WL}}$ (\%) & $\Delta_{\text{WLM-WL}}$ (\%) & $\Delta_{\text{RAW-WL}}$ (\%) \\
\midrule
swe & 0.0099 & 0.0099 & --      & 0.0096 & 0.0  & --   & -3.0 \\
dan & 0.1099 & 0.1122 & 0.1393  & 0.1215 & +2.1 & +26.8 & +10.6 \\
eng & 0.1256 & 0.1566 & 0.2593  & 0.2113 & +24.7 & +106.5 & +68.2 \\
spa & 0.1094 & 0.1419 & 0.1641  & 0.1453 & +29.7 & +50.0 & +32.8 \\
lat & 0.1136 & 0.1329 & 0.1595  & 0.1422 & +17.0 & +40.4 & +25.2 \\
fre & 0.1687 & 0.2485 & 0.2983  & 0.2945 & +47.3 & +76.8 & +74.5 \\
ita & 0.1050 & 0.1196 & 0.1553  & 0.1553 & +13.8 & +47.9 & +47.9 \\
\midrule
\textbf{Average} & \textbf{0.1060} & \textbf{0.1316} & \textbf{0.1960} & \textbf{0.1548}
& \textbf{+24.2} & \textbf{+85.0} & \textbf{+46.1} \\
\bottomrule
\end{tabular}
\end{table*}

% TODO: Add results for byT5 or the chosen prompted model.

% TODO: Add results for ipa2ipa|from=X|to=sv and comparison against prompt-only transfer.

\section{Discussion}

% TODO: Discuss the baseline as motivation: simple multilingual mixing and labels do not solve the intended foreign-word-in-Swedish-phoneset task.
% TODO: Discuss whether the pivot task reduces the amount of direct donor-language-to-Swedish-target data required.
% TODO: Discuss whether explicit labels are an acceptable assumption, given the SSML/text-processing use case.

\section{Conclusion}

% TODO: Write only after the prompted and pivot experiments are run.

\section{Acknowledgements}

Place all acknowledgments (including those concerning research grants and funding) in a separate section at the end of the paper.

\nocite{*}
\section{Bibliographical References}\label{sec:reference}

\bibliographystyle{lrec2026-natbib}
\bibliography{lrec2026-example}

\section{Language Resource References}
\label{lr:ref}
\bibliographystylelanguageresource{lrec2026-natbib}
\bibliographylanguageresource{languageresource}

\end{document}

%%% Local Variables:
%%% mode: latex
%%% TeX-master: t
%%% End:
